\section{Công trình liên quan}
\label{sec:related-work}

\subsection{Retrieval-Augmented Generation và hỏi--đáp nhiều bước}

Retrieval-Augmented Generation (RAG) được đề xuất như một kiến trúc kết hợp
khả năng sinh của mô hình ngôn ngữ với một bộ nhớ tri thức bên ngoài
\cite{lewis2020rag}. Thay vì chỉ dựa vào tri thức được mã hóa trong tham số
mô hình, RAG truy hồi các đoạn văn liên quan đến truy vấn và sử dụng chúng
làm ngữ cảnh cho quá trình sinh. Cơ chế này giúp cập nhật nguồn tri thức mà
không cần huấn luyện lại toàn bộ mô hình, đồng thời tạo điều kiện để truy
vết câu trả lời về các tài liệu nguồn.

Phần lớn pipeline RAG cơ bản thực hiện truy hồi ở mức đoạn văn và xếp hạng
các ứng viên dựa trên độ tương đồng giữa biểu diễn truy vấn và biểu diễn tài
liệu. Cách tiếp cận này giả định rằng một hoặc một số đoạn được xếp hạng cao
chứa đủ thông tin để trả lời câu hỏi. Tuy nhiên, giả định đó không luôn đúng
đối với câu hỏi nhiều bước. Trong những trường hợp này, bằng chứng thứ nhất
có thể xác định một sự kiện trung gian, còn bằng chứng thứ hai mới liên kết
sự kiện trung gian với kết luận cuối cùng. Nếu mỗi bằng chứng được xếp hạng
độc lập, hệ thống có thể truy hồi đúng từng mảnh thông tin nhưng không nhận
ra thứ tự suy luận giữa chúng.

HotpotQA là một trong những benchmark có ảnh hưởng đối với hỏi--đáp nhiều
bước, trong đó hệ thống không chỉ phải đưa ra câu trả lời mà còn phải xác
định các supporting facts được sử dụng
\cite{yang2018hotpotqa}. Các nghiên cứu trên dạng tác vụ này cho thấy chất
lượng truy hồi từng đoạn riêng lẻ và chất lượng khôi phục toàn bộ chuỗi suy
luận là hai vấn đề khác nhau. Điều này đặc biệt quan trọng trong miền pháp
luật, nơi thứ tự áp dụng quy tắc có thể quyết định ý nghĩa của kết luận:
một điều kiện ban đầu làm phát sinh trạng thái pháp lý trung gian và trạng
thái đó tiếp tục kích hoạt một hệ quả pháp lý khác.

Bên cạnh cấu trúc truy hồi, chất lượng biểu diễn truy vấn và tài liệu cũng
ảnh hưởng trực tiếp đến khả năng tìm đúng rule hoặc event. BGE-M3 hỗ trợ
nhiều cơ chế biểu diễn và truy hồi trong một mô hình đa ngôn ngữ, bao gồm
dense retrieval, sparse retrieval và multi-vector retrieval
\cite{chen2024bgem3}. Trong nghiên cứu này, BGE-M3 được sử dụng để mã hóa
các rule và event tiếng Việt. Tuy nhiên, embedding chỉ cung cấp tín hiệu
tương đồng ngữ nghĩa; cấu trúc suy luận nhiều bước vẫn cần được khôi phục
thông qua đồ thị và quá trình mở rộng causal path.

\subsection{Graph-based RAG và truy hồi theo đường dẫn}

Graph-based RAG thay thế hoặc bổ sung kho đoạn văn độc lập bằng một biểu
diễn đồ thị, trong đó node mô tả thực thể, sự kiện, khái niệm hoặc tài liệu,
còn edge mô tả quan hệ giữa các node. Cấu trúc này cho phép truy hồi không
chỉ dựa trên độ tương đồng trực tiếp với truy vấn mà còn dựa trên các node
lân cận, đường đi và cộng đồng trong graph.

GraphRAG tổ chức tri thức được trích xuất thành graph và sử dụng các bản tóm
tắt theo cộng đồng để hỗ trợ truy vấn cần tổng hợp thông tin ở phạm vi rộng
\cite{edge2024graphrag}. Mục tiêu chính của cách tiếp cận này là khai thác
cấu trúc toàn cục của corpus thay vì chỉ lấy các đoạn văn gần nhất với truy
vấn. HippoRAG khai thác graph như một dạng bộ nhớ liên kết, cho phép lan
truyền tín hiệu từ các khái niệm được kích hoạt bởi truy vấn tới những node
có quan hệ nhiều bước \cite{gutierrez2024hipporag}. Các nghiên cứu này cho
thấy graph có thể cải thiện khả năng kết nối những bằng chứng không đồng
thời có độ tương đồng cao với câu hỏi.

Dù vậy, graph connectivity không tự động tương đương với tính hợp lệ của
một chuỗi suy luận. Một đường đi có thể tồn tại vì hai sự kiện cùng xuất
hiện trong một tài liệu, cùng liên quan đến một thực thể hoặc được kết nối
bởi một loại quan hệ không có hướng nhân quả. Trong miền pháp luật, sự khác
biệt này càng quan trọng vì hệ thống cần phân biệt quan hệ tham chiếu, quan
hệ cùng điều luật và quan hệ điều kiện--hệ quả.

Nghiên cứu này sử dụng một graph dị thể trong đó causal path được giới hạn
bởi các cạnh \texttt{CAUSES} hình thành từ quy tắc pháp luật đã chuẩn hóa.
Mỗi hop của đường dẫn phải được liên kết với ít nhất một rule và căn cứ điều
luật tương ứng. Do đó, đơn vị truy hồi chính không chỉ là node hoặc đoạn văn
mà là một chuỗi có thứ tự:

\[
v_0 \xrightarrow{r_1} v_1
    \xrightarrow{r_2} \cdots
    \xrightarrow{r_h} v_h,
\]

trong đó \(v_i\) là sự kiện và \(r_i\) là quy tắc pháp luật hỗ trợ bước
chuyển. Biểu diễn này không khẳng định rằng mọi cạnh trong graph là một quan
hệ nhân quả thực nghiệm; nó chỉ mã hóa hướng suy diễn quy phạm từ điều kiện
đến hệ quả theo corpus được sử dụng.

\subsection{Mô hình nhân quả cấu trúc và suy luận phản thực tế}

Structural Causal Model cung cấp ngôn ngữ hình thức để biểu diễn các biến,
cơ chế cấu trúc và kết quả của can thiệp
\cite{pearl2009causality}. Điểm quan trọng phân biệt can thiệp với quan sát
là phép toán \(do(\cdot)\). Khi thực hiện \(do(M=m)\), giá trị của biến
\(M\) được gán cưỡng bức và cơ chế thông thường tạo ra \(M\) bị thay thế.
Sau đó, các biến hậu duệ được suy luận lại theo những cơ chế cấu trúc còn
lại. Cách diễn giải này khác với việc chỉ xóa một node khỏi graph và kiểm
tra xem giữa hai node khác còn tồn tại đường đi hay không.

Các nghiên cứu gần đây đã bắt đầu kết hợp causal reasoning với RAG.
CausalRAG khai thác cấu trúc nhân quả để hỗ trợ quá trình truy hồi và suy
luận thay vì chỉ dựa trên liên kết ngữ nghĩa
\cite{wang2025causalrag}. Các hướng causal-counterfactual RAG tiếp tục khảo
sát cách sử dụng claim phản thực tế hoặc thay đổi ngữ cảnh để kiểm tra tính
ổn định của câu trả lời
\cite{khadilkar2025causalcounterfactualrag}. CF-RAG thể hiện sự quan tâm
ngày càng tăng đối với counterfactual reasoning trong các hệ thống RAG
\cite{qin2026cfrag}. Các công trình này cho thấy việc bổ sung cấu trúc nhân
quả có thể hỗ trợ phân biệt giữa bằng chứng chỉ có tương quan ngữ nghĩa và
bằng chứng có vai trò trong một cơ chế suy luận.

Nghiên cứu hiện tại có phạm vi khác với causal discovery hoặc causal effect
estimation. Hệ thống không học hướng cạnh từ dữ liệu quan sát và không ước
lượng một đại lượng tác động trung bình. Thay vào đó, hướng điều
kiện--hệ quả được lấy từ các quy tắc pháp luật đã chuẩn hóa. LegalSCM sử
dụng các rule này như cơ chế tất định để kiểm tra factual path và mô phỏng
việc gán một mediator về trạng thái \(\mathrm{FALSE}\). Vì vậy, kết quả
counterfactual của LegalSCM chỉ có ý nghĩa trong phạm vi mô hình quy phạm và
corpus đầu vào.

Một khác biệt khác của nghiên cứu này là phân tích claim theo truy vấn. Hệ
thống không thực hiện can thiệp trên mọi câu hỏi. Đối với câu hỏi nhân quả
thông thường, LegalSCM chủ yếu kiểm tra khả năng kích hoạt factual path. Đối
với claim về cạnh trực tiếp, hệ thống kiểm tra cấu trúc edge thay vì suy
diễn rằng một đường hai bước tương đương với quan hệ trực tiếp. Hard
intervention chỉ được diễn giải khi truy vấn thực sự đề cập đến việc loại
bỏ hoặc giả định sự vắng mặt của một mediator.

\subsection{Hỏi--đáp và suy luận trong miền pháp luật}

Các benchmark pháp luật hiện có bao phủ nhiều dạng bài toán, từ phân loại
và dự đoán kết quả đến truy hồi căn cứ và trả lời câu hỏi. CaseHOLD đánh giá
khả năng lựa chọn holding phù hợp trong các vụ án, qua đó nhấn mạnh yêu cầu
phân biệt giữa những phương án pháp lý có nội dung gần nhau
\cite{zheng2021casehold}. LegalBench tập hợp nhiều tác vụ nhằm đánh giá các
năng lực pháp lý khác nhau của mô hình ngôn ngữ, bao gồm nhận diện quy tắc,
áp dụng quy tắc và suy luận trên tình huống
\cite{guha2023legalbench}.

Trong khi đó, LegalBench-RAG tập trung rõ hơn vào vai trò của retrieval
trong các ứng dụng pháp luật \cite{alami2024legalbenchrags}. Hướng nghiên
cứu này phù hợp với yêu cầu thực tế rằng câu trả lời pháp lý cần được gắn
với tài liệu hoặc điều luật nguồn. Tuy nhiên, việc truy hồi đúng tài liệu
chưa bảo đảm rằng mô hình đã áp dụng các quy tắc theo đúng thứ tự. Citation
cũng có thể đúng ở mức điều luật nhưng chưa chắc hỗ trợ trực tiếp kết luận
được sinh ra.

So với các benchmark trên, bộ đánh giá được sử dụng trong nghiên cứu này có
phạm vi hẹp hơn về hệ thống pháp luật nhưng tập trung cụ thể vào cấu trúc
suy luận đa bước trong Bộ luật Hình sự Việt Nam. Mỗi câu hỏi được liên kết
với gold rule, gold event, gold path, gold answer và gold citation. Các dạng
câu hỏi bao gồm suy luận thuận, suy luận ngược, tìm bridge event, giải thích
chuỗi, kiểm tra cạnh trực tiếp, phân nhánh và hội tụ.

Tuy nhiên, bộ đánh giá hiện tại được sinh từ cùng graph nguồn và chưa trải
qua quy trình thẩm định đầy đủ bởi chuyên gia pháp luật. Do đó, nghiên cứu
không coi đây là một benchmark pháp lý có tính chuẩn tắc hoặc có khả năng
đại diện cho mọi dạng suy luận trong Bộ luật Hình sự. Nó được sử dụng như
một testbed có kiểm soát để đánh giá khả năng khôi phục đường dẫn, xác minh
claim và sinh câu trả lời từ graph.

\subsection{Định vị nghiên cứu}

Nghiên cứu này không xem từng thành phần riêng lẻ---embedding retrieval,
graph traversal, SCM hoặc extractive answer generation---là một đóng góp
hoàn toàn mới. Đóng góp nằm ở cách các thành phần được tích hợp và giới hạn
cho bài toán suy luận pháp luật:

\begin{enumerate}
    \item Rule và event được biểu diễn đồng thời trong bộ nhớ vector và đồ
    thị nhân quả dị thể.

    \item Kết quả retrieval được tổ chức thành causal path có thứ tự, trong
    đó mỗi hop được gắn với rule và căn cứ pháp luật.

    \item Query analysis quyết định loại claim và dạng câu trả lời mà không
    sử dụng nhãn benchmark.

    \item LegalSCM kiểm tra factual path và hard intervention trong phạm vi
    mô hình quy phạm, trong khi node-deletion reachability được giữ như một
    baseline riêng biệt.

    \item Câu trả lời cuối cùng được sinh theo answer intent và chỉ sử dụng
    citation từ những rule thực sự tham gia vào đường suy luận được chọn.

    \item Structural SCM và path ablation được đánh giá theo paired setup để
    phân biệt tác động lên chất lượng với chi phí tính toán và khả năng diễn
    giải.
\end{enumerate}

So sánh thực nghiệm cho thấy structural SCM và path ablation đạt các quality
metric giống nhau trên benchmark hiện tại. Kết quả này khác với một tuyên
bố về tính vượt trội của SCM: nó chỉ ra rằng tập câu hỏi hiện có chưa đủ
nhạy để phân biệt hai semantics can thiệp. Phát hiện này đồng thời đặt ra
yêu cầu xây dựng các mẫu có ground truth về factual world, mediator
intervention và counterfactual outcome trong các nghiên cứu tiếp theo.
